\capitulo{1}{Introducción}

La Esclerosis Lateral Amiotrófica (ELA) es una enfermedad neurodegenerativa progresiva que afecta selectivamente a las neuronas motoras, destruyendo de forma irreversible el control muscular voluntario y comprometiendo de manera drástica, en fases avanzadas, la comunicación oral de las personas que la padecen \cite{startstemcells:als}. A medida que la enfermedad progresa, más del 80\% de los pacientes experimentan una pérdida gradual de la inteligibilidad del habla debido al deterioro bulbar, evolucionando desde la disartria hasta la anartria completa \cite{cave2021voice}. Esta pérdida del canal verbal dificulta gravemente la interacción con su entorno habitual y genera un impacto directo en su autonomía personal, su autoidentidad y su bienestar emocional.

En este contexto, los Sistemas de Comunicación Aumentativa y Alternativa (SAAC) desempeñan un papel fundamental \cite{plenainclusion:saac}, ya que permiten a las personas con ELA mantener la capacidad de comunicarse mediante diferentes tecnologías adaptadas a sus limitaciones y al estadio de la enfermedad \cite{pmc2025stage}. Actualmente existe una amplia variedad de SAAC, que difieren en cuanto a su nivel tecnológico, métodos de acceso (interfaces visuales, pulsadores o control cefálico) y plataformas disponibles, lo que convierte el proceso de selección del sistema óptimo en una tarea altamente compleja tanto para los pacientes como para sus familias y los profesionales sanitarios.

Este trabajo se enmarca dentro del ámbito de la Ingeniería de la Salud y aborda esta problemática desde una perspectiva de apoyo a la toma de decisiones. Aunque ya existen  catálogos oficiales y guías clínicas estáticas que recopilan información sobre los distintos SAAC disponibles \cite{ceapat:productos, elaandalucia:guia}, estos suelen presentarse en formatos extensos y estáticos que requieren tiempo, conocimientos previos y esfuerzo por parte del usuario para extraer conclusiones. En la práctica clínica real, especialmente en el caso de enfermedades de progresión acelerada como la ELA, la ventana temporal para intervenir de manera preventiva es crítica, por lo que los métodos de consulta estáticos y manuales resultan ineficientes \cite{taylor2026timing}.

La motivación de este trabajo surge, por tanto, de la necesidad de transformar la información técnica existente en una herramienta más accesible, estructurada y orientada a la toma de decisiones. Desde la Ingeniería de la Salud es posible formalizar el conocimiento disponible sobre los SAAC y convertirlo en criterios y reglas que faciliten la selección del sistema más adecuado según las necesidades de cada usuario. Además, hará posible reconocer aquellos casos en los que los sistemas existentes no cubren completamente dichas necesidades, facilitará la identificación de carencias y  permitirá desarrollar propuestas de mejora desde un enfoque de diseño inclusivo.

El objetivo principal de este trabajo es realizar una revisión y comparación técnica de los SAAC utilizados por personas con ELA y, a partir de dicho análisis, diseñar una herramienta de apoyo a la decisión que permita orientar sobre las opciones más adecuadas en función de cada perfil. 

La memoria se estructura de la siguiente manera:
El capítulo 2 presenta los objetivos específicos, técnicos y formativos del trabajo. 
El capítulo 3 desarrolla el marco teórico, revisando la literatura existente sobre la ELA, los sistemas SAAC disponibles y las herramientas de apoyo a la decisión en este ámbito. 
El capítulo 4 describe la metodología seguida durante el desarrollo, incluyendo la arquitectura de la aplicación, el diseño de la base de datos y el algoritmo de filtrado.
El capítulo 5 presenta los resultados obtenidos tras la implementación y validación del sistema. Incluye también los enlaces correspondientes a la página desarrollada y el repositorio GitHub. 
El capítulo 6 recoge la discusión de dichos resultados, analizando su alcance y limitaciones. 
El capítulo 7 expone las conclusiones del trabajo. 
Finalmente, el capítulo 8 propone las líneas de trabajo futuras que podrían continuarse a partir de este proyecto.
